\section{Theory}\label{sec:theory}
% =============================================================================================
% Owner: theory writer. Source: docs/theory.md (Sections 1-9), checked against src/synthval.
% Proofs: sections/appendix_proofs.tex (input by appendix.tex under \section{Proofs}).
% Labels follow paper/OUTLINE.md, Section 6. New labels (not in the outline registry):
%   sec:theory-markowitz, ass:anc, ass:rel, prop:anchored-target, rem:anchored-null,
%   rem:nonstat-past, rem:ks-marginal, tab:theory-status, app:proofs-* (appendix subsections),
%   cor:block-bias, rem:calibrated-bad, tab:manual-theory-calibration
%   (paper/scripts/theory_calibration.py), tab:manual-e1-sim-stat (paper/scripts/e1_sim_stat.py).
% =============================================================================================

This section states what the protocol of \cref{sec:method} guarantees and what it does not. A
Proposition, Lemma or Corollary is proved, in \cref{app:proofs} unless the proof takes a few
lines; ``sketch'' marks a proof that gives every key step but omits routine measure-theoretic
detail. A Fact is a standard result quoted with a reference. A Remark whose title contains
``heuristic'' is an argument, not a theorem. \Cref{tab:theory-status} at the end of the section
lists the status and the assumptions of each result. Throughout, $P$ and $Q$ are the laws of one
real and one synthetic window on $\R^d$ ($d=L$, or $d=NL$ for multivariate windows), $\norm{\cdot}$
is the Euclidean (Frobenius) norm, $Pf=\int f\,dP$, and $\Pev f$, $\Qev f$, $\Ptr f$, $\Qtr f$
are means of $f$ over the evaluation and training halves of the two samples.

% ---------------------------------------------------------------------------------------------
\subsection{What a window test can detect}\label{sec:theory-windows}

Let $(r_t)_{t\in\mathbb Z}$ be a return process with values in $\R$ or $\R^N$. The real sample
consists of non-overlapping windows $X_i=(r_{s_0+(i-1)L+1},\dots,r_{s_0+iL})$ of the observed
series, and the synthetic sample of windows $Y_j$ of paths simulated from the generator $G$.

\begin{definition}[Absolute null]\label{def:abs-null}
With $P=\Law(X_i)$ and $Q=\Law(Y_j)$, the absolute null is $H_0:P=Q$, against $H_1:P\neq Q$.
\end{definition}

$H_0$ concerns the marginal law of one window. Dependence between windows, such as successive
windows of one path or windows of different tickers on the same dates, is not part of $H_0$; it is
a nuisance that changes the null distribution of every statistic.
\Cref{sec:theory-critic,sec:theory-level,sec:theory-consistency,sec:theory-likelihood} assume
independent windows and \cref{sec:theory-dependence} removes this assumption.

\begin{proposition}[What a window law identifies]\label{prop:window-ident}
Let $(r_t)$ be strictly stationary and let $P_L$ be the law of $(r_1,\dots,r_L)$.
\begin{enumerate}
  \item[(a)] $P_L$ determines the joint law of $(r_{t_1},\dots,r_{t_k})$ whenever
  $t_k-t_1\le L-1$: every autocovariance and every moment $\E g(r_t,r_{t+h})$ with
  $0\le h\le L-1$ (the autocorrelations of $r_t$, $\abs{r_t}$ and $r_t^2$, leverage correlations),
  and the law of aggregated returns $\sum_{t=1}^h r_t$ for $h\le L$.
  \item[(b)] The family $(P_L)_{L\ge1}$ determines the law of the process.
  \item[(c)] For every $L\ge1$ there are strictly stationary processes with the same $P_L$ and
  different $P_{L+1}$.
\end{enumerate}
\end{proposition}
\noindent\textit{Proof.} In \cref{app:proofs-windows}.

A test on windows of length $L$ therefore tests the necessary condition $P_L=Q_L$: equality of
every dependence feature at lags $0,\dots,L-1$ and of aggregated returns up to horizon $L$. It has
no power against deviations that appear only at lags $L$ or more, however large the sample.
Longer windows see longer-range dependence but give fewer windows ($n=\lfloor T/L\rfloor$) in a
higher dimension; E5 varies $L\in\{5,10,20,60\}$ (\cref{sec:robust}).

\begin{proposition}[Pairs determine a stationary Markov chain]\label{prop:markov-pairs}
Let $(r_t)$ and $(r'_t)$ be strictly stationary, time-homogeneous, first-order Markov chains with
values in a Polish space. If $(r_0,r_1)\eqd(r'_0,r'_1)$, then $(r_t)\eqd(r'_t)$ as processes. For
chains of order $k$, $P_{k+1}=P'_{k+1}$ implies equality of the processes.
\end{proposition}
\noindent\textit{Proof.} In \cref{app:proofs-windows}.

\Cref{prop:markov-pairs} compares two Markov chains: if only the generator is Markov, $P_2=Q_2$
says that the real process has the generator's pair law, nothing more. The models of this paper
are Markov only in an augmented state. GBM, Merton and the i.i.d.\ Student-$t$ laws of E2 have
i.i.d.\ returns, and a window density is a product of daily densities. In the \GARCHt{} model,
$\sigma_{t+1}^2=\omega/(1-b)+a\sum_{j\ge0}b^j\varepsilon_{t-j}^2$ is a function of past returns
(the model is observation-driven), so the likelihood is explicit given $\sigma_1^2$; yet the
returns are not Markov of any finite order when $a,b>0$ (sketch in \cref{app:proofs-windows}).
The conditional variance has persistence $a+b$: shocks to $\sigma_t^2$ decay like $(a+b)^h$, and
so does the autocorrelation of the squared shocks $\varepsilon_t^2$ when $\E\varepsilon_t^4<\infty$
\citep{francq_2019_garch}. With
standardised Student-$t_\nu$ innovations that requires $(a+b)^2+2a^2(\nu-1)/(\nu-4)<1$, which fails
for the E2 parameters ($a=0.08$, $b=0.90$, $\nu=5$), so for them only the statement about
$\sigma_t^2$ applies. At persistence $0.98$ the volatility half-life is
$\log0.5/\log0.98\approx34$ days, beyond $L=20$.
In the Heston model the variance is a latent continuous state with its own shock and no closed-form
filter, so the likelihood is intractable \citep{aitsahalia2007maximum}; we simulate it with the
quadratic-exponential scheme of \citet{andersen_2008_simple} (\cref{sec:data-generators}). The
stationary bootstrap has no likelihood. Likelihood benchmarks thus exist for GBM, Merton and
\GARCHt{} only, and the critic is the only formal tool of the paper for the other two generators.

\begin{remark}[Pooling tickers]\label{rem:pooling}
Real and synthetic returns of ticker $i$ are divided by the same constant $s_i$ computed on
\TRAIN. Given \TRAIN{} this is a fixed bijection, so the per-ticker null $P_i=Q_i$ is unchanged.
If ticker $i$ contributes a fraction $\pi_i$ of the real windows and $\pi'_i$ of the synthetic
ones, the pooled laws are $\bar P=\sum_i\pi_iP_i$ and $\bar Q=\sum_i\pi'_iQ_i$, and the per-ticker
nulls imply the pooled null $\bar P=\bar Q$ only if $\pi_i=\pi'_i$. Each ticker therefore
contributes the same number of real and synthetic windows. The pooled null is weaker than the
intersection of the per-ticker nulls, because per-ticker errors can cancel in the mixture; the
per-ticker tests of \cref{rem:bh} address this.
\end{remark}

% ---------------------------------------------------------------------------------------------
\subsection{The held-out critic statistic}\label{sec:theory-critic}

Each sample is split into a training and an evaluation half, with index sets $\Itr,\Iev$ (real,
sizes $n_t,n_e$) and $\Jtr,\Jev$ (synthetic, sizes $m_t,m_e$). Let $\Gsig$ be the $\sigma$-field
generated by the training halves, the index sets and the training randomness $\xi$
(initialisation, mini-batches). The critic is $\fhat=\ghat\circ A$, where
$A(x)=\Lambda^{-1}(x-m_A)$ is the coordinate-wise standardisation fitted on the pooled training
half (per asset channel for multivariate windows) and $\ghat$ is the network trained by
\eqref{eq:wgan-gp}. Conditioning on $\Gsig$ fixes $\fhat$, and the statistic is
$T=\Pev\fhat-\Qev\fhat$ \eqref{eq:critic-stat}. The networks are compositions of globally
Lipschitz maps (affine and convolutional layers, LeakyReLU, LayerNorm with a positive
$\epsilon$), so $\Lip(\ghat)<\infty$ and $\fhat\in L^p(P)$ whenever $P\in\Pspace{p}(\R^d)$, although
the constant is not controlled.

\begin{fact}[Kantorovich--Rubinstein duality]\label{fact:kr-duality}
\citep[Ch.~5--6]{villani2009optimal}. For $P,Q\in\Pspace{1}(\R^d)$,
$\W(P,Q):=\inf_{\pi\in\Pi(P,Q)}\int\norm{x-y}\,\pi(dx,dy)=\sup_{f\in\Lip_1}(Pf-Qf)$, where
$\Pi(P,Q)$ is the set of couplings and $\Lip_c=\{f:\Lip(f)\le c\}$. $\W$ is a metric on
$\Pspace{1}$, and it metrises weak convergence together with convergence of first moments.
\end{fact}

$\W$ is the integral probability metric of $\Lip_1$ \citep{muller1997integral}; on $\Pspace{1}$,
$H_0$ holds if and only if $\W(P,Q)=0$.

\begin{remark}[The critic sees standardised windows]\label{rem:representation}
Given $\Gsig$, $A$ is a fixed affine bijection, so $A_\#P=A_\#Q$ if and only if $P=Q$. The
gradient penalty differentiates $\ghat$ with respect to the standardised input, so the Lipschitz
constant it pushes towards one is $\Lip(\ghat)$, the Lipschitz constant of $\fhat$ for the metric
$\norm{x-y}_A=\norm{\Lambda^{-1}(x-y)}$. Write $P_A=A_\#P$ and $Q_A=A_\#Q$. Then $\W(P_A,Q_A)$ is
the Wasserstein-1 distance of $P$ and $Q$ under $\norm{\cdot}_A$, and
$\norm{\Lambda}_{\mathrm{op}}^{-1}\W(P,Q)\le\W(P_A,Q_A)\le\norm{\Lambda^{-1}}_{\mathrm{op}}\W(P,Q)$.
In E2 the windows are raw daily returns with volatility $0.02$, so $\norm{\Lambda^{-1}}_{\mathrm{op}}$
is about $50$; in E3 and E4 the returns are already divided by $s_i$ and $\Lambda$ is of order one.
\end{remark}

\begin{proposition}[Conditional unbiasedness and a bound by $\W$ of the standardised laws]
\label{prop:unbiased}
Assume that, conditionally on $\Gsig$, every evaluation window $X_i$ has law $P$, every
evaluation window $Y_j$ has law $Q$, and $P,Q\in\Pspace{1}(\R^d)$.
\begin{enumerate}
  \item[(a)] $\E[T\mid\Gsig]=P\fhat-Q\fhat=P_A\ghat-Q_A\ghat=:\Dcrit(\fhat)$.
  \item[(b)] $\Dcrit(\fhat)\le\Lip_{S_A}(\ghat)\,\W(P_A,Q_A)\le\Lip(\ghat)\,\W(P_A,Q_A)$, where
  $\Lip_{S_A}$ is the Lipschitz constant on $S_A=\supp P_A\cup\supp Q_A$.
  \item[(c)] If the evaluation windows are also conditionally independent and
  $P,Q\in\Pspace{2}(\R^d)$, then
  $\Var(T\mid\Gsig)=\Var_P\fhat/n_e+\Var_Q\fhat/m_e\le\Lip(\ghat)^2\big(\Tr\Sigma_{P_A}/n_e+\Tr\Sigma_{Q_A}/m_e\big)$,
  where $\Sigma_{P_A}=\Lambda^{-1}\Sigma_P\Lambda^{-1}$ is the covariance matrix of $P_A$.
\end{enumerate}
\end{proposition}
\noindent\textit{Proof.} In \cref{app:proofs-critic}.

\begin{remark}[What the $\W$ reading of $T$ does and does not say]\label{rem:gp-soft}
If $\Lip(\ghat)\le1$, part (b) makes the conditional mean of $T$ a lower bound on $\W(P_A,Q_A)$;
this is the sense in which $T$ is ``a $\W$ lower bound''. The reading has five limits.
(i) \emph{Scale:} it concerns standardised windows, not return units (\cref{rem:representation}).
(ii) \emph{Sampling error:} $T$ itself can exceed $\W(P_A,Q_A)$ by a term of order
$n_e^{-1/2}+m_e^{-1/2}$; the conservative summary is the lower confidence bound
$T-z_{1-\alpha}(s_X^2/n_e+s_Y^2/m_e)^{1/2}$.
(iii) \emph{Soft constraint:} the penalty enforces the Lipschitz constraint only in expectation
over interpolates of training points, and it is two-sided, pushing gradient norms towards one
rather than below one \citep{petzka_2018_regularization}. Nothing guarantees $\Lip(\ghat)\le1$,
and WGAN critics are known to be poor estimators of $\W$ itself
\citep{mallasto2019well,stanczuk_2021_wasserstein}.
(iv) \emph{Looseness:} a small $T$ says that the trained critic found no separating direction, not
that $\W$ is small; \cref{prop:consistency} says what training would have to deliver.
(v) \emph{Relevance:} $\W$ of standardised windows weights directions by the inverse coordinate
scale and ignores what a user of the data values; \cref{sec:theory-markowitz} shows that it can
rank generators against their portfolio utility.
We therefore use $T$ as a test statistic. Its level (\cref{sec:theory-level}) does not depend on
any Lipschitz property, and its $\W$ reading is approximate and descriptive.
\end{remark}

\begin{remark}[Why hold data out]\label{rem:holdout}
The plug-in distance $\W(P_n,Q_m)$ is biased upwards, since $\E\,\W(P_n,P)$ decays at the rate
$n^{-1/d}$ for $d\ge3$ \citep{fournier2015rate}, with a constant that depends on $d$ and on $P$. With
one window the bias is $\E\,\W(\delta_{X_1},P)=\E\norm{X_1-X_1'}$, $X_1'$ an independent copy, which
is of the order of the scale of the windows. The rate factor is still $0.68$ at $d=20$ and
$n=2000$, so at this rate going from one window to 2000 removes only about a third of that bias.
The in-sample training objective is biased upwards in the same way. The
held-out $T$ is conditionally unbiased for $\Dcrit(\fhat)$ and has an exact null distribution
(\cref{prop:exact-level}), at the cost of sample splitting, which \cref{prop:local-power}
quantifies. This is the template of sample-split two-sample tests
\citep{lopezpaz2017revisiting,kim2021classification,liu2020learning,kubler2022witness}; for
Wasserstein-based testing see \citet{ramdas2017wasserstein}. The critic is a learned, sample-split
relative of the test functions of \citet{bierens1990consistent} and
\citet{stinchcombe1998consistent}.
\end{remark}

% ---------------------------------------------------------------------------------------------
\subsection{Exact level under exchangeability}\label{sec:theory-level}

Pool the evaluation windows as $V=(V_1,\dots,V_{n_e+m_e})=(X_{\Iev},Y_{\Jev})$, real first, and let
$t(v)=n_e^{-1}\sum_{k\le n_e}\fhat(v_k)-m_e^{-1}\sum_{k>n_e}\fhat(v_k)$. Then $T=t(V)$, the
permuted statistics are $T_b=t(\pi_b\cdot V)$ with $(\pi\cdot v)_k=v_{\pi^{-1}(k)}$, and the
p-value \eqref{eq:perm-pvalue} is $\pperm=(1+\#\{b\ge1:T_b\ge T\})/(1+B)$.

\begin{assumption}[Exchangeability]\label{ass:exch}
\emph{(A1)} (i) $X_1,\dots,X_n$ are i.i.d.\ $P$, $Y_1,\dots,Y_m$ are i.i.d.\ $Q$, and the two
samples are independent. (ii) The index sets are chosen independently of the window values;
random, deterministic, group-wise and time-block splits qualify. (iii) $\fhat$ is a measurable
function of the training windows and of randomness $\xi$ independent of the data. (iv) The
permutations $\pi_1,\dots,\pi_B$ are i.i.d.\ uniform on the permutations of $\{1,\dots,n_e+m_e\}$,
independent of everything else.\\
\emph{(A1$'$)} Conditionally on $\Gsig$, the law of $V$ is invariant under every permutation of
its coordinates. Under $H_0$, (A1) implies (A1$'$).
\end{assumption}

\begin{proposition}[Exact level]\label{prop:exact-level}
Under (A1$'$) and (A1)(iv), in particular under (A1) and $H_0$, for every $\alpha\in(0,1)$, every
$B\ge1$, every $n_e,m_e\ge1$ and every critic, however it was trained,
\[
  \Prob(\pperm\le\alpha\mid\Gsig)\le\frac{\lfloor\alpha(B+1)\rfloor}{B+1}\le\alpha\quad\text{a.s.},
  \qquad\text{hence}\qquad \Prob(\pperm\le\alpha)\le\alpha .
\]
\end{proposition}
\noindent\textit{Proof.} In \cref{app:proofs-level} (\cref{lem:exchangeable,lem:rank}).

This is the randomisation principle \citep{lehmann2022testing} in the random-permutation form of
\citet{phipson2010permutation} and \citet{hemerik2018exact}. The guarantee needs no Lipschitz
constraint, no convergence of training and no continuity: a poor or overfitted critic costs power,
never level. For a discrete statistic such as the C2ST accuracy, randomised tie-breaking makes the
p-value exactly uniform under (A1$'$). The guarantee does need conditional exchangeability of the
evaluation windows given the training data, which fails when real windows depend on each other or
on the training windows (\cref{sec:theory-dependence}).

\begin{remark}[Weaker nulls and studentisation]\label{rem:studentised}
(a) Conditionally on $\Gsig$, $z=T/(s_X^2/n_e+s_Y^2/m_e)^{1/2}$ is asymptotically $N(0,1)$
whenever $P\fhat=Q\fhat$ and second moments are finite (Lindeberg); under $H_0$ its normal p-value
is asymptotically equivalent to $\pperm$ \citep{hoeffding_1952_large}.
(b) The permutation test is exact for $P=Q$. For the weaker hypothesis $H_0':P\fhat=Q\fhat$ the
unstudentised permutation test can be miscalibrated when $n_e\neq m_e$ and the two variances
differ \citep{romano_1990_behavior}; with $n_e=m_e$ the permutation variance matches the true
variance asymptotically. Permuting the studentised statistic is asymptotically valid for $H_0'$
and stays exact under $H_0$ \citep{janssen1997studentized,chung2013exact}. These results concern
unequal i.i.d.\ laws, not serial dependence. They matter for the relative tests, whose nulls are
statements about means of a trained critic.
\end{remark}

\begin{remark}[Seeds and the p-value lottery]\label{rem:combine-seeds}
The p-value depends on the random split, on $\xi$ and, on real data, on the synthetic seed.
Reporting the smallest p-value over $S$ seeds is invalid. If every $p_s$ is valid, that is
$\Prob(p_s\le u)\le u$ for all $u\in(0,1)$, then $\pmed=\min\{1,2\med_sp_s\}$ \eqref{eq:pmed} is
valid under arbitrary dependence between the $p_s$ \citep{meinshausen_2009_pvalues}: the event
$\{\pmed\le\alpha\}$ implies that at least $S/2$ of the $p_s$ are at most $\alpha/2$, and Markov's
inequality bounds its probability by $(2/S)\sum_s\Prob(p_s\le\alpha/2)\le\alpha$. Twice the mean
is valid as well \citep{vovk_2020_combining}. The combination inherits any miscalibration of the
individual $p_s$.
\end{remark}

% ---------------------------------------------------------------------------------------------
\subsection{Consistency}\label{sec:theory-consistency}

The standardisation $A$ is estimated on the training half. Under the conditions of
\cref{ass:consistency} below, if moreover the real share $n_t/(n_t+m_t)$ of the training windows
converges and the limiting pooled law has positive coordinate variances, $A$ converges in
probability (almost surely for nested training halves) to a population standardisation $A^*$, and
the estimation error adds $o_\Prob(1)$ terms to every step (\cref{app:proofs-consistency}).
We therefore state the results for the standardised laws $A^*_\#P$ and $A^*_\#Q$ and for the
network $\ghat$, with Lipschitz constants in standardised coordinates, and, in this subsection
only, write $P,Q$ for $A^*_\#P,A^*_\#Q$ and $\fhat$ for $\ghat$.

\begin{assumption}[Consistency regime]\label{ass:consistency}
\emph{(C1)} (A1) holds, $P,Q\in\Pspace{2}(\R^d)$ and $P\neq Q$.
\emph{(C2)} $n_t,m_t,n_e,m_e\to\infty$ and $B\ge\lceil2/\alpha\rceil$ ($\pperm\le\alpha$ requires
only $B+1\ge1/\alpha$; the stronger condition is used in the proof).
\emph{(C3)} (Hard Lipschitz constraint, an idealisation of the gradient penalty.) There are
$c<\infty$ and classes $\mathcal F_n\subset\Lip_c$ with $\Prob(\fhat\in\mathcal F_n)\to1$.
\emph{(C4)} (Richness.) There are $n_0$ and $f_0\in\bigcap_{n\ge n_0}\mathcal F_n$ with
$Pf_0-Qf_0=:\eta_0>0$.
\emph{(C5)} (Approximate optimisation.) With probability tending to one,
$(\Ptr-\Qtr)\fhat\ge\sup_{f\in\mathcal F_n}(\Ptr-\Qtr)f-\varepsilon_{\mathrm{opt}}$ with
$\varepsilon_{\mathrm{opt}}\le\eta_0/2$.
\end{assumption}

\begin{proposition}[Consistency of the critic test]\label{prop:consistency}
Under \cref{ass:consistency}, $\Prob(\pperm\le\alpha)\to1$.
\end{proposition}
\noindent\textit{Proof.} In \cref{app:proofs-consistency} (uses \cref{lem:uniform-lip}).

\begin{corollary}[$T$ estimates $\W$ of the standardised laws under idealised training]
\label{cor:w1-limit}
Suppose (C1)--(C3) hold with $c=1$, $\sup_{f\in\mathcal F_n}(Pf-Qf)\to\W(P,Q)$, and (C5) holds
with $\varepsilon_{\mathrm{opt}}\to0$ in probability. Then $T\pto\W(A^*_\#P,A^*_\#Q)$, the
Wasserstein-1 distance under the metric $\norm{\Lambda_*^{-1}(x-y)}$, not in return units.
\end{corollary}
\noindent\textit{Proof.} In \cref{app:proofs-consistency}.

\begin{remark}[How realistic are (C3)--(C5)?]\label{rem:c-realism}
(C4) is a population property. It holds if, for some $c$, the $c$-Lipschitz members of
$\bigcup_n\mathcal F_n$ are dense in the bounded 1-Lipschitz functions for uniform convergence on
compacts, because bounded Lipschitz functions separate probability measures
\citep{dudley_2002_real}; norm-constrained GroupSort networks have this density property
\citep{anil_2019_sorting}. For the LeakyReLU and LayerNorm networks used here, trained with a soft
penalty, we do not prove it. (C3) idealises the penalty (\cref{rem:gp-soft}), and (C5) is an
assumption about stochastic-gradient training that we cannot verify. \Cref{prop:consistency} says
what has to hold for consistency, not that WGAN-GP training delivers it; the power study of
\cref{sec:sim-power} measures what training delivers.
\end{remark}

% ---------------------------------------------------------------------------------------------
\subsection{Likelihood connection and the Neyman--Pearson bound}\label{sec:theory-likelihood}

\begin{proposition}[Bayes-optimal logit; \citealp{goodfellow2014generative}]\label{prop:bayes-logit}
Let $\nu=P+Q$, $p=dP/d\nu$, $q=dQ/d\nu$, and let $(V,\iota)$ satisfy
$\iota\sim\mathrm{Bernoulli}(\pi)$, $V\mid\iota=1\sim P$ and $V\mid\iota=0\sim Q$. Among measurable
$D:\R^d\to[0,1]$, the cross-entropy risk $-\E[\iota\log D(V)+(1-\iota)\log(1-D(V))]$ is minimised
by $D^*=\pi p/(\pi p+(1-\pi)q)$, $\nu$-a.e., whose logit is
$\lstar(v)=\log\{p(v)/q(v)\}+\log\{\pi/(1-\pi)\}$.
\end{proposition}
\begin{proof}
The risk equals $\int-[\pi p\log D+(1-\pi)q\log(1-D)]\,d\nu$. Where $p,q>0$ the integrand is
strictly convex in $D\in[0,1]$ and minimised at $D^*$; where $q=0<p$ the minimiser is $D=1$, and
conversely.
\end{proof}

\begin{corollary}[Targets of the logit and accuracy statistics]\label{cor:jeffreys}
(a) If $\lhat=\lstar+\text{const}$ and $\KL(P\Vert Q)$, $\KL(Q\Vert P)$ are finite, then the
\LRC{} statistic $\Tlogit=\Pev\lhat-\Qev\lhat$ satisfies
$\E[\Tlogit\mid\Gsig]=\KL(P\Vert Q)+\KL(Q\Vert P)=\Jeff(P,Q)$, the Jeffreys divergence; the class
prior cancels. (b) For a decision region $\hat A=\{\lhat>0\}$ the balanced accuracy satisfies
$\E[\Tacc\mid\Gsig]=\tfrac12+\tfrac12\{P(\hat A)-Q(\hat A)\}\le\tfrac12\{1+\TV(P,Q)\}$, with
equality for $\hat A=\{p>q\}$.
\end{corollary}
\begin{proof}
Apply \cref{prop:unbiased}(a), which needs no Lipschitz property, to $f=\lstar$ and to
$f=\one_{\hat A}$.
\end{proof}

Unlike the \WC{}, nothing makes $\E[\Tlogit\mid\Gsig]$ a lower bound on $\Jeff$ for an imperfect
logit: rescaling $\lhat$ inflates it. $\Tlogit$ is a plug-in estimator, consistent only if
$\lhat\to\lstar$ in $L^1(P)\cap L^1(Q)$; the level of the test does not depend on this. The
\LRC{} is thus a learned log-likelihood-ratio statistic \citep{cheng2022classification}, and
\cref{sec:sim-llr} measures how close the learned logit is to $\log(p/q)$ (\cref{tab:e2-llr}).

\begin{proposition}[No two-sample test beats the Neyman--Pearson oracle]\label{prop:oracle-bound}
Let $P,Q$ have densities $p,q$ with respect to a common $\sigma$-finite measure. Write
$\bNP(F\to F';m,\alpha)$ for the power of the most powerful level-$\alpha$ test of
$F^{\otimes m}$ against $F'^{\otimes m}$ \citep{lehmann2022testing}. Let
$\phi(X_{1:n},Y_{1:m},U)\in[0,1]$ be any test whose auxiliary randomness $U$ (splits, training,
permutations) is independent of the data.
\begin{enumerate}
  \item[(a)] If $\E_{P^{\otimes n}\otimes P^{\otimes m}}\phi\le\alpha$, then
  $\E_{P^{\otimes n}\otimes Q^{\otimes m}}\phi\le\bNP(P\to Q;m,\alpha)$.
  \item[(b)] If $\E_{Q^{\otimes n}\otimes Q^{\otimes m}}\phi\le\alpha$, then
  $\E_{P^{\otimes n}\otimes Q^{\otimes m}}\phi\le\bNP(Q\to P;n,\alpha)$.
\end{enumerate}
\end{proposition}
\noindent\textit{Proof.} In \cref{app:proofs-likelihood}.

A two-sample test of exact level $\alpha$ thus has power at most
$\min\{\bNP(P\to Q;m,\alpha),\bNP(Q\to P;n,\alpha)\}$: its level must hold at every null point
$P=Q$, while the oracle knows both densities. The E2 oracle \eqref{eq:lr-oracle} applies the
Neyman--Pearson statistic $\TLR=\sum_j\log\{q(Y_j)/p(Y_j)\}$ to the synthetic sample with a Monte
Carlo critical value. As a Monte Carlo test it has exact level, and its power tends to $\bNP$ as
the number of null samples grows \citep{dufour_2006_monte}. For the i.i.d.\ families (a), (b) and
(d) of E2 the window density is a product of daily densities and the oracle is exact. For the
\GARCHt{} family (c) the window density
$p(x_{1:L})=\int\prod_tp(x_t\mid\sigma_t^2)\,\Pi(d\sigma_1^2)$, with $\Pi$ the law of the
conditional variance at the window start, has no closed form. E2 replaces it by an average over
$M=512$ draws from $\Pi$, shared by the observed and the null samples. The Monte Carlo test keeps
its exact level, but its statistic only approximates the log-likelihood ratio (the logarithm of an
unbiased density estimate is biased downwards by a term of order $1/M$). Its power need not equal
$\bNP$, so the family-(c) curve is an approximate envelope, and a learned test above it would not
contradict \cref{prop:oracle-bound}.

\begin{proposition}[Local power of fixed-critic tests]\label{prop:local-power}
Let $\{Q_\theta\}_{\theta\in[0,\theta_0)}$, with $Q_0=P$, be differentiable in quadratic mean at
$0$ with score $h\in L^2(P)$ and $I=Ph^2>0$. Let $m\to\infty$ with $n/m\to\kappa\in(0,\infty)$,
$n_e=\lfloor\rho n\rfloor$ and $m_e=\lfloor\rho m\rfloor$ for some $\rho\in(0,1]$, and
$\theta_m=\tau/\sqrt m$ with $\tau>0$.
\begin{enumerate}
  \item[(a)] $\bNP(P\to Q_{\theta_m};m,\alpha)\to1-\Phi(z_{1-\alpha}-\tau\sqrt I)$.
  \item[(b)] Let $f\in L^2(P)$ be fixed, with $\sigma_f^2=\Var_Pf>0$ and orientation
  $r_f=-\Corr_P(f,h)\ge0$ (the critic scores real windows higher). The test that rejects when
  $z>z_{1-\alpha}$, and the permutation test with the same $f$ and $B\to\infty$, have power
  under
  $P^{\otimes n}\otimes Q_{\theta_m}^{\otimes m}$ tending to
  $1-\Phi\big(z_{1-\alpha}-\tau\sqrt I\,r_f\sqrt{\rho\kappa/(1+\kappa)}\big)$.
  \item[(c)] At $\kappa=1$ and $\rho=1/2$ (the defaults), the noncentrality of every test based on
  the held-out mean difference of a fixed critic (studentised or permutation) is at most half the
  oracle's. Where the oracle has power $0.80$ at $\alpha=0.05$, every such test has asymptotic
  power at most $0.34$.
\end{enumerate}
\end{proposition}
\noindent\textit{Proof.} In \cref{app:proofs-likelihood}; proved for a fixed critic, sketched for
a learned one.

The factor $\sqrt{\rho\kappa/(1+\kappa)}$ separates the two costs of not knowing the densities:
the two-sample cost $\sqrt{\kappa/(1+\kappa)}$, which tends to one only when $P$ is effectively
known, and the splitting cost $\sqrt\rho$. The factor $r_f\le1$ is the price of a sub-optimal
direction. The bound concerns the mean-difference statistic $T$ that the protocol uses; other
statistics of the held-out critic values, such as rank or quantile comparisons, are not covered.
With a learned critic the argument applies conditionally on the training half with $f=\fhat$, so
the bound in (c) holds for mean-difference tests with learned critics too; a rigorous version
needs a central limit theorem uniform over the critic class, which we do not provide.

\begin{example}[Volatility level]\label{ex:vol-level}
Let $P=N(0,\sigma^2I_L)$ and $Q_\theta=N(0,(1+\theta)^2\sigma^2I_L)$. The score is proportional to
$\norm{x}^2-L\sigma^2$. The critic $\pm\norm{x}$, oriented to score real windows higher, is an
optimal Kantorovich potential for every scale factor, and its local efficiency is
$r_f^2=\Corr(\chi_L,\chi_L^2)^2=0.876$, $0.957$, $0.988$ for $L=1,5,20$. The population
standardisation is a multiple of the identity here, so standardising changes nothing. Against
volatility-level deviations the population \WC{} is nearly as efficient as the likelihood-ratio
direction (details in \cref{app:proofs-likelihood}).
\end{example}

\begin{remark}[\LRC{} versus \WC{} (heuristic)]\label{rem:lr-vs-w1}
By \cref{prop:bayes-logit} the population logit critic is $\log(p/q_\theta)\approx-\theta h$, the
locally optimal direction ($r_f=1$). The population \WC{} is a Kantorovich potential, which is in
general not proportional to $h$ ($r_f<1$). If both were learned well, the \LRC{} would be locally
at least as powerful. The Lipschitz penalty may instead reduce the estimation variance of the
\WC{} in samples of a few thousand windows. Which effect dominates is an empirical question; in
E2 the \WC{} is the more powerful of the two in most cells (\cref{sec:sim-power}).
\end{remark}

\begin{remark}[What the pooled KS test can and cannot see]\label{rem:ks-marginal}
(i) Suppose the returns are i.i.d.\ within windows under both laws, $P=P_1^{\otimes L}$ and
$Q=Q_1^{\otimes L}$, as in E2 families (a) and (b). Then the pair formed by the empirical
distribution of the $nL$ pooled real returns and that of the $mL$ pooled synthetic returns is
sufficient for $(P_1,Q_1)$. For any test $\phi$, the conditional expectation of $\phi$ given this
pair has the same rejection probability at every $(P_1,Q_1)$, so nothing is lost by looking only
at the two pooled marginals. This is what KS does, with no splitting and no learning; the oracle
statistic $\TLR=\sum_{j,t}\log\{q_1(y_{jt})/p_1(y_{jt})\}$ is also a function of the pooled
synthetic marginal. A window critic must learn this additive structure in dimension $L$ from half
of the data. This is the regime in which E2 finds KS far ahead of every window test
(\cref{tab:e2-power}).
(ii) When returns within a window are dependent, the null distribution of KS, derived for $nL$
independent draws, is wrong: in E2 family (c) KS is oversized under $H_0$ (\cref{tab:e2-size}).
(iii) When $P$ and $Q$ have the same one-dimensional marginal and differ only in their dependence
(the E2-dep families), the scaled KS statistic converges in law to the supremum of a Gaussian
process under the alternative as under the null, so its rejection probability converges to a
limit below one: KS is not consistent against dependence-only deviations, whereas the window tests
are consistent under \cref{ass:consistency} if the deviation is visible at lags below $L$
(\cref{prop:window-ident}; \cref{fig:e2dep-power}). Parts (ii) and (iii) follow from the
empirical-process central limit theorem for window-level sums of indicators; we do not write out
the details.
\end{remark}

\begin{remark}[Likelihood ranks by $\KL$, the critic by $\W$]\label{rem:kl-vs-w1}
The expected out-of-sample log-likelihood of a model $G$ is $\E_P\log g=-\mathrm{Ent}(P)-\KL(P\Vert G)$
whenever finite, and for conditional densities of stationary ergodic processes it is the
corresponding relative-entropy rate. Rankings by predictive log-score, and the Vuong and
Amisano--Giacomini tests \citep{vuong1989likelihood,amisano2007comparing}, therefore rank by
$\KL(P\Vert G)$; the critic ranks by (lower bounds on) $\W$ of window laws. Both vanish only at
$G=P$, but they can rank misspecified models differently. Let $P$ be Student-$t_4$ scaled to unit
variance, $G_1=N(0,1)$ (right variance, wrong tails) and $G_2$ the same $t_4$ law scaled by $1.4$
(right tails, wrong scale). Numerical quadrature gives $\KL(P\Vert G_1)\approx0.084$,
$\W(P,G_1)\approx0.139$, $\KL(P\Vert G_2)\approx0.060$ and $\W(P,G_2)=0.4\,\E_P\abs{X}\approx0.283$:
the likelihood prefers $G_2$, $\W$ prefers $G_1$. Disagreement between the two evaluations in
\cref{sec:emp-likelihood} therefore shows that they measure different divergences, not that
either is broken.
\end{remark}

% ---------------------------------------------------------------------------------------------
\subsection{The relative tests}\label{sec:theory-relative}

There are three samples: $\Rnew$, windows $X_1,\dots,X_n$ with law $\Pnew$; the synthetic sample,
$Y_1,\dots,Y_m$ with law $\QG$; and $\Rpast$, windows $Z_1,\dots,Z_k$ with law $\Ppast$. Markets
are non-stationary, and a generator calibrated on the past inherits features of the past, such
as its volatility level, that need not persist; on our data the absolute \WC{} test rejects every
generator (\cref{tab:e3-pooled-tests}), whereas the window tests do not reject the past windows
themselves against the new ones (E1(c), \cref{tab:e1-nonstat}), a comparison for which they have
little power with about 15 \TEST{} evaluation dates (\cref{sec:emp-placebo}). The relative question is whether $G$ is at
least as close to $\Rnew$ as $\Rpast$ is. For a sample $W$ and a part $\mathrm{p}\in\{\trn,\evl\}$
of it, $W^{\mathrm p}f$ denotes the mean of $f$ over that part.

\emph{Anchored test} (primary, \cref{sec:method-relative}). $Y$ and $Z$ are split into training
and evaluation parts of sizes $m_t,m_e$ and $k_t,k_e$; all of $X$ is used for training. One critic
$\fhat$ is trained to separate $X$ (positive class) from the pool $\{Y^{\trn},Z^{\trn}\}$, in which
the two samples carry weight $1/2$ each. The statistic and its conditional target are
\[
  \Danc=Z^{\evl}\fhat-Y^{\evl}\fhat\quad\eqref{eq:anchored-stat},\qquad
  \delta=\Ppast\fhat-\QG\fhat ,
\]
where now $\Gsig$ is generated by $X$, $Y^{\trn}$, $Z^{\trn}$, the index sets and $\xi$. Large
$\Danc$ means that the critic which scores ``$\Rnew$-likeness'' finds the past more $\Rnew$-like
than $G$.

\emph{Two-critic test} (secondary). Critics $\fhat_G$ on $(X^{\trn},Y^{\trn})$ and $\fhat_P$ on
$(X^{\trn},Z^{\trn})$ are trained separately with independent randomness and evaluated on the
shared $X^{\evl}$: $\Dtwo=(X^{\evl}-Y^{\evl})\fhat_G-(X^{\evl}-Z^{\evl})\fhat_P$
\eqref{eq:two-critic-stat}, with targets $D_G=(\Pnew-\QG)\fhat_G$, $D_P=(\Pnew-\Ppast)\fhat_P$ and
$\delta^{\mathrm{2c}}=D_G-D_P$.

\begin{definition}[Relative nulls]\label{def:rel-null}
(a) The anchored null is $H_0^{\mathrm{anc}}:\delta\le0$, that is $\QG\fhat\ge\Ppast\fhat$ for the
trained critic, against $\delta>0$. (b) The two-critic null is
$H_0^{\mathrm{2c}}:\delta^{\mathrm{2c}}\le0$. (c) The idealised $\W$ null is
$H_0^{W}:\W(\QG,\Pnew)\le\W(\Ppast,\Pnew)$. The boundary point $\QG=\Ppast$ belongs to $H_0^W$ and
to the analogous null for any divergence.
\end{definition}

Both conditional nulls depend on the data through the trained critics. \Cref{prop:anchored}
gives exact or asymptotic level for $H_0^{\mathrm{anc}}$; \cref{prop:anchored-target} says what
$\delta$ measures; \cref{prop:paired-boot,prop:boundary-size} treat the two-critic test.

\begin{assumption}[Anchored exchangeability]\label{ass:anc}
(i) Conditionally on $X$, the $k+m$ windows $Z_1,\dots,Z_k,Y_1,\dots,Y_m$ are i.i.d. This holds
when the three samples are independent, each is i.i.d., and $\QG=\Ppast$, whatever $\Pnew$.
(ii) The split of $Y$ and $Z$ depends only on time stamps, group labels and randomness independent
of the window values. (iii) $\fhat$ is a measurable function of $X$, $Y^{\trn}$, $Z^{\trn}$ and of
randomness $\xi$ independent of the data. (iv) The permutations of the $k_e+m_e$ evaluation labels
are i.i.d.\ uniform and independent of everything else.
\end{assumption}

\begin{proposition}[Level of the anchored test]\label{prop:anchored}
\begin{enumerate}
  \item[(a)] \emph{Exact level at the boundary.} Under \cref{ass:anc}, the permutation p-value of
  $\Danc$ (labels of the pooled $Z^{\evl}$ and $Y^{\evl}$ permuted, $\fhat$ fixed) satisfies
  $\Prob(\pperm\le\alpha\mid\Gsig)\le\lfloor\alpha(B+1)\rfloor/(B+1)\le\alpha$ almost surely, for
  every critic, every $B\ge1$ and $k_e,m_e\ge1$, and whatever $\Pnew$.
  \item[(b)] \emph{Weaker null, independent windows.} Suppose that, conditionally on $\Gsig$,
  $Z^{\evl}$ is i.i.d.\ $\Ppast$, $Y^{\evl}$ is i.i.d.\ $\QG$, the two are independent, and
  $\fhat$ has finite positive variances under both laws. As $k_e,m_e\to\infty$ with
  $k_e/(k_e+m_e)\to\kappa\in(0,1)$ and $\fhat$ fixed, the studentised permutation test satisfies
  $\limsup\Prob(\text{reject}\mid\Gsig)\le\alpha$ on $\{\delta\le0\}$. When $k_e=m_e$, so does the
  unstudentised permutation test.
  \item[(c)] \emph{Dependent windows.} Suppose that, conditionally on $\Gsig$, the evaluation
  windows of $Z$ and of $Y$, ordered by date rank, are independent of each other and satisfy
  \cref{ass:mixing} with marginals $\Ppast$ and $\QG$ and $f=\fhat$. Then the Driscoll--Kraay
  statistic $\zhac$ of $\Danc$ \eqref{eq:hac-z} is asymptotically $N(0,1)$ on $\{\delta=0\}$ and
  tends to $-\infty$ in probability on $\{\delta<0\}$, so the test that rejects when
  $\zhac>z_{1-\alpha}$ has asymptotic level $\alpha$ for $H_0^{\mathrm{anc}}$.
\end{enumerate}
\end{proposition}
\noindent\textit{Proof.} In \cref{app:proofs-relative}.

\begin{proposition}[What the anchored target measures]\label{prop:anchored-target}
Lipschitz constants and $\W$ are taken in the standardised coordinates $A$ of the anchored critic,
and $M=\tfrac12(\QG+\Ppast)$ is the law of the pool.
\begin{enumerate}
  \item[(a)] If $\Ppast,\QG\in\Pspace{1}$ and $\Lip(\ghat)\le c$, then
  $\abs{\delta}\le c\,\W(A_\#\Ppast,A_\#\QG)$. In particular $\delta=0$ whenever $\QG=\Ppast$,
  whatever $\Pnew$ and the critic.
  \item[(b)] \emph{(\LRC{}.)} Let $\Pnew,\QG,\Ppast$ be mutually absolutely continuous with
  $\KL(\QG\Vert\Pnew)$ and $\KL(\Ppast\Vert\Pnew)$ finite. If $\fhat=\lstar$ is the Bayes logit of
  the population training problem ($\Pnew$ against $M$ with equal class weights), then
  \begin{equation}\label{eq:anchored-kl}
    \delta=\big[\KL(\QG\Vert\Pnew)-\KL(\Ppast\Vert\Pnew)\big]-\big[\KL(\QG\Vert M)-\KL(\Ppast\Vert M)\big],
  \end{equation}
  and the second bracket lies in $[-\log2,\log2]$. If moreover $d\QG/d\Ppast=1+\varepsilon h$ with
  $h$ bounded and $\Ppast h=0$, the second bracket equals $-\varepsilon^3\Ppast(h^3)/24+O(\varepsilon^4)$
  as $\varepsilon\to0$.
  \item[(c)] \emph{(\WC{}.)} If $\fhat=f^*$ is a 1-Lipschitz maximiser of $(\Pnew-M)f$, then
  $\delta=(\Pnew-\QG)f^*-(\Pnew-\Ppast)f^*$, and the two terms sum to $2\,\W(A_\#\Pnew,A_\#M)$.
\end{enumerate}
\end{proposition}
\noindent\textit{Proof.} In \cref{app:proofs-relative}.

\begin{remark}[What the anchored null means away from the boundary]\label{rem:anchored-null}
By \cref{prop:anchored-target}(a) the anchored test is a two-sample test of $G$ against the past
whose witness has been trained to point towards $\Rnew$: it can reject only if $G$ differs from the
past, and only in a direction that separates $\Rnew$ from the pool. If the market were stationary
($\Pnew=\Ppast$) it would be a critic test of $\QG=\Ppast$. By part (b), near the boundary the
idealised \LRC{} version compares $\KL(\QG\Vert\Pnew)$ with $\KL(\Ppast\Vert\Pnew)$ up to a
third-order term, so to first order $H_0^{\mathrm{anc}}$ says that $G$ is at least as close to the
new law as the past is, in the Kullback--Leibler divergence from the candidate to the new law; this
is the reverse of the direction $\KL(\Pnew\Vert G)$ by which predictive log-scores rank
(\cref{rem:kl-vs-w1}). Away from the boundary the correction in \eqref{eq:anchored-kl} is bounded
but not small. By part (c), the \WC{} version compares the two absolute discrepancies with a common
witness. Each is at most the corresponding $\W$ distance, but a common witness need not be optimal
for either pair, so $\delta>0$ is possible under $H_0^W$ and conversely: neither relative test is a
test of $H_0^W$. A rejection (\BAD) means that $G$ is detectably less $\Rnew$-like than the past
for a critic trained to recognise $\Rnew$ against an even mixture of both. A non-rejection is not
evidence that $G$ is good, which is why the verdict needs a power check (\cref{sec:theory-verdict}).
Two idealisations remain. First, $G$ is calibrated on $\Rpast$, so its windows depend on the past
windows through the fitted parameters and \cref{ass:anc}(i) holds only approximately; for a
parametric generator each past window moves the estimate by $O(1/k)$, which we expect to be
negligible but do not prove. Second, the stationary bootstrap copies past windows that the critic
saw as pool negatives; evaluation windows of $G$ then duplicate training points, exchangeability
fails in the direction of rejection, and the generator receives the verdict \INVALID{} (E2 part e
documents this leakage, \cref{tab:e2-relative-boundary}).
\end{remark}

\begin{remark}[Non-stationary past (heuristic)]\label{rem:nonstat-past}
On real data $\Rpast$ is the \TRAIN{} panel, whose windows are not identically distributed (the
2020 crash, the 2022 drawdown). \Cref{ass:anc}(i) then fails, and the natural boundary is
$\QG$ equal to the pooled \TRAIN{} window law. The time-block split sends a random half of the
\TRAIN{} blocks to the evaluation part, so $Z^{\evl}\fhat$ estimates the critic mean over a random
half of the \TRAIN{} regimes while $\QG\fhat$ averages over all of them. The target $\delta$ is then
a random variable, not necessarily centred at zero, whose spread reflects the heterogeneity between
blocks, which the Driscoll--Kraay variance of the evaluation part does not capture. By the argument of
\cref{prop:boundary-size}(a), the rejection rate at this boundary is about
$\E[1-\Phi(z_{1-\alpha}-\delta/\sigma)]$, which exceeds $\alpha$ when $\delta$ is roughly symmetric
and not small relative to the evaluation standard deviation $\sigma$. In addition, the last
\TRAIN{} block borders the first \TEST{} dates; with volatility clustering, when that block falls
in the evaluation part it raises $Z^{\evl}\fhat$, in the direction of rejection. We therefore claim
no calibration for a single-seed anchored p-value on real data. The verdict uses $\pmed$ over the
synthetic seeds, and the size of this rule at the pooled-\TRAIN{} boundary and its power are
estimated in design D2, which keeps the real \TRAIN{} volatility path, as well as in the
stationary design D1 (\cref{def:power-check}, \cref{tab:e3-power}). The estimate depends on the
periods: at the main split the boundary size of the rule in D2 is 0.08 [0.04, 0.15], at the
alternative split of E5 it is 0.60 [0.39, 0.79] (Clopper--Pearson 95\% intervals;
\cref{tab:manual-e5-split}). \Cref{rem:calibrated-bad} states when the boundary replications
calibrate an observed $\pmed$, and separates the power of the rule from its miscalibration.
\end{remark}

For the two-critic test we use the regime of \citet{bounliphone2016test}-type comparisons with
learned witnesses.

\begin{assumption}[Two-critic regime]\label{ass:rel}
(R1) Conditionally on $\Gsig$ (the training halves of the three samples and the training
randomness of both critics), the three evaluation samples are mutually independent and i.i.d.\
from $\Pnew$, $\QG$ and $\Ppast$. (R2) $n_e,m_e,k_e\to\infty$ proportionally. With
$\sigma_N^2=\sigma_X^2/n_e+\sigma_Y^2/m_e+\sigma_Z^2/k_e$, $\sigma_X^2=\Var_{\Pnew}(\fhat_G-\fhat_P)$,
$\sigma_Y^2=\Var_{\QG}\fhat_G$ and $\sigma_Z^2=\Var_{\Ppast}\fhat_P$, the quantity
$(n_e+m_e+k_e)\sigma_N^2$ is bounded away from zero and the $(2+\epsilon)$-th moments of the
three summands are bounded in probability. (R3) The bootstrap resamples $X^{\evl}$ once per
replicate for both critics (paired), and $Y^{\evl}$, $Z^{\evl}$ independently.
\end{assumption}

\begin{proposition}[Paired bootstrap, conditional validity; sketch]\label{prop:paired-boot}
Under \cref{ass:rel}, conditionally on $\Gsig$: (a) $(\Dtwo-\delta^{\mathrm{2c}})/\sigma_N\dto N(0,1)$;
(b) the bootstrap law of $\Dtwo{}^*-\Dtwo$ estimates the law of $\Dtwo-\delta^{\mathrm{2c}}$
uniformly consistently; (c) the test that rejects when
$p^*=(1+\#\{b:\Dtwo_b{}^*-\Dtwo\ge\Dtwo\})/(1+B)\le\alpha$ has
$\Prob(\text{reject}\mid\Gsig)=1-\Phi(z_{1-\alpha}-\delta^{\mathrm{2c}}/\sigma_N)+o_\Prob(1)$ as
$B\to\infty$, which is at most $\alpha+o_\Prob(1)$ on $\{\delta^{\mathrm{2c}}\le0\}$; (d) the
percentile interval has asymptotic coverage $0.95$ for $\delta^{\mathrm{2c}}$.
\end{proposition}
\noindent\textit{Proof.} Sketch in \cref{app:proofs-relative}. The centred p-value follows
\citet{hall1991two}.

\begin{proposition}[The two-critic test is anti-conservative at the boundary]
\label{prop:boundary-size}
Assume \cref{ass:rel} along a sequence of experiments, let $B\to\infty$, $\alpha<1/2$ and
$z=z_{1-\alpha}$.
\begin{enumerate}
  \item[(a)] $\Prob(\text{reject})=\E[1-\Phi(z-\delta^{\mathrm{2c}}/\sigma_N)]+o(1)$.
  \item[(b)] Suppose in addition that $\QG=\Ppast$, $m_t=k_t$ and $m_e=k_e$; that
  $\fhat_G=\Psi(X^{\trn},Y^{\trn},\xi_G)$ and $\fhat_P=\Psi(X^{\trn},Z^{\trn},\xi_P)$ for one
  training map $\Psi$ with $\xi_G,\xi_P$ i.i.d.\ and independent of the data; and that the joint
  law of the data and the index sets is invariant under exchanging $(Y,\xi_G)$ and $(Z,\xi_P)$.
  Then $(\delta^{\mathrm{2c}},\sigma_N)\eqd(-\delta^{\mathrm{2c}},\sigma_N)$ and
  \[
    \Prob(\text{reject})=\E\,\bar g\big(\abs{\delta^{\mathrm{2c}}}/\sigma_N\big)+o(1),\qquad
    \bar g(x)=1-\tfrac12\{\Phi(z-x)+\Phi(z+x)\}.
  \]
  $\bar g$ is continuous and strictly increasing on $[0,\infty)$, with $\bar g(0)=\alpha$ and
  $\bar g(\infty)=1/2$. Hence $\Prob(\text{reject})\ge\alpha+o(1)$, and if
  $\Prob(\abs{\delta^{\mathrm{2c}}}>\varepsilon\sigma_N)\ge\eta$ for some $\varepsilon,\eta>0$ and
  all large samples, then $\liminf\Prob(\text{reject})\ge\alpha+\eta\{\bar g(\varepsilon)-\alpha\}>\alpha$.
  \item[(c)] If the training halves are held fixed while the evaluation sizes grow, then
  $\Prob(\text{reject})\to\Prob(\delta^{\mathrm{2c}}>0)+\alpha\Prob(\delta^{\mathrm{2c}}=0)$. Under
  the conditions of (b) this limit equals $1/2$ when $\delta^{\mathrm{2c}}$ has no atom at zero.
\end{enumerate}
\end{proposition}
\noindent\textit{Proof.} In \cref{app:proofs-relative}.

At the boundary, $\delta^{\mathrm{2c}}$ is the difference of two exchangeable replicates of a
training-dependent quantity. It fluctuates at the scale of critic training, and nothing in
WGAN-GP training makes it $o_\Prob(\sigma_N)$. The proof uses only the conditional rejection
probability and the symmetry, so the conclusion also holds for the HAC version of the two-critic
test. The anchored test avoids the problem: with one critic the target is exactly zero at the
boundary (\cref{prop:anchored-target}(a)) and the permutation is exact (\cref{prop:anchored}(a)).
E2 part e compares the two tests at the boundary (\cref{sec:sim-relative}). With $n=500$ windows
per sample the two-critic test rejects at 0.040 and 0.065 at nominal 5\% in its two boundary
designs, and the anchored test at 0.065 and 0.055 (\cref{tab:e2-relative-boundary}): in that
design the spread of $\delta^{\mathrm{2c}}$ induced by training is small relative to $\sigma_N$,
and the excess allowed by part (b) is not visible. The anchored test is preferred because its level
is exact for every critic and every sample size, not because the two-critic test failed there.

\begin{remark}[Guarantee of the two-critic test; the MMD version]\label{rem:rel-guarantee}
Under \cref{ass:rel}, a two-critic rejection means that the critic family, given equal resources,
separates $G$ from $\Rnew$ by more than it separates $\Rpast$ from $\Rnew$. By
\cref{prop:unbiased}, $D_G$ and $D_P$ equal $\Lip(\ghat_G)\W(\QG,\Pnew)$ and
$\Lip(\ghat_P)\W(\Ppast,\Pnew)$ minus unknown non-negative slacks (in standardised coordinates),
so with respect to $H_0^W$ the test is neither exact nor conservative. It is fair only to the
extent that the two critic problems are equally hard (same architecture, budget and sample sizes),
and even then \cref{prop:boundary-size} shows over-rejection at $\QG=\Ppast$. The MMD version of
\citet{bounliphone2016test} uses a fixed kernel and a genuine metric null, but its normal
approximation needs a non-degenerate projection, which fails when all three laws coincide. For
serially dependent evaluation windows the i.i.d.\ bootstrap of $X^{\evl}$ must be replaced by a
block bootstrap \citep{kunsch1989jackknife,politis1994stationary,lahiri2003resampling} or by a HAC
variance, as implemented.
\end{remark}

% ---------------------------------------------------------------------------------------------
\subsection{Serial and cross-sectional dependence}\label{sec:theory-dependence}

The real windows of one ticker form a stationary but dependent sequence: volatility clustering
makes window features such as realised variance persistent. Windows of different tickers on the
same date are dependent through common factors, whereas synthetic windows from univariate
generators are independent across tickers. Under $H_0$ the pooled evaluation vector is then not
exchangeable and \cref{prop:exact-level} does not apply.

\begin{remark}[Size distortion of the naive permutation test (heuristic)]\label{rem:naive-perm}
Let $a_k=\fhat(X_{(k)})$ be the real evaluation values in time order, with variance $\sigma^2$ and
long-run variance $\LRV_a=\sigma^2(1+2\sum_{h\ge1}\rho_h)$, and let the synthetic values be
i.i.d.\ with variance $\sigma^2$, with $n_e=m_e$. The permutation distribution uses the variance
$2\sigma^2/n_e$, while the true variance is $(\LRV_a+\sigma^2)/n_e$, so the rejection rate tends
to $1-\Phi\big(z_{1-\alpha}\{2/(2+2\sum_h\rho_h)\}^{1/2}\big)$. For $\rho_h=\phi^h$ the nominal 5\%
test rejects about 12\% of the time at $\phi=0.5$ and 17\% at $\phi=0.67$. Contemporaneous
covariances between tickers inflate the true variance in the same way.
\end{remark}

For a fixed function $f$, let $a_k=f(X_k)$ be the real evaluation values, each with the date rank
$\tau(k)$ of its first day, and let $\mathcal S_\tau=\sum_{k:\tau(k)=\tau}(a_k-\bar a)$ be the
date sums. The Driscoll--Kraay variance \eqref{eq:dk-var} is the Bartlett-kernel long-run variance
of $(\mathcal S_\tau)$, with lags counted in windows and the truncation lag $H$ chosen by the
\citet{andrews1991heteroskedasticity} AR(1) plug-in; with one window per date it is the
\citet{newey1987simple} estimator. The synthetic values $b_k$ are treated in the same way, and
$\zhac=(\bar a-\bar b)/\{\hat V_{\mathrm{DK}}(\bar a)+\hat V_{\mathrm{DK}}(\bar b)\}^{1/2}$
\eqref{eq:hac-z}.

\begin{assumption}[Weak dependence, fixed critic $f$]\label{ass:mixing}
Each sample has a fixed number of windows per date ($N$ for a panel, one for a single series).
(D1) The date-level sums $u_\tau=\sum_{k:\tau(k)=\tau}f(X_k)$ of the real sample and
$v_\tau=\sum_{k:\tau(k)=\tau}f(Y_k)$ of the synthetic sample form strictly stationary sequences,
independent of each other, that are strongly mixing with coefficients satisfying
$\sum_h\alpha(h)^{\epsilon/(2+\epsilon)}<\infty$, and $u_\tau$, $v_\tau$ have finite
$(2+\epsilon)$-th moments for some $\epsilon>0$ (for Lipschitz $f$ this follows from finite
$(2+\epsilon)$-th moments of returns). (D2) The long-run variances $\LRV_u=\sum_h\Cov(u_0,u_h)$
and $\LRV_v$ are positive. (D3) The Driscoll--Kraay estimators are consistent:
$\hat V_{\mathrm{DK}}(\bar a)/\Var(\bar a)\pto1$, and likewise for $\bar b$; sufficient conditions,
including data-dependent lags, are in \citet{andrews1991heteroskedasticity},
\citet{newey1987simple}, \citet{driscoll1998consistent} and \citet{dejong_2000_consistency}.
(D4) The numbers of evaluation dates of the two samples tend to infinity proportionally.
\end{assumption}

With the time-block split the evaluation dates of a sample form a union of at most $K$ contiguous
blocks; because $K$ is fixed while the blocks grow, (D1) for the underlying date sequence carries
the central limit theorem and (D3) over to the union. The edge terms are of order $KH^2$ divided by
the number of dates, which vanishes when $H$ grows like the cube root of that number, the
\citet{andrews1991heteroskedasticity} rate (\cref{app:proofs-dependence}).

\begin{proposition}[Asymptotic validity of the HAC test for a fixed critic]\label{prop:hac}
Under \cref{ass:mixing} and $H_0':Pf=Qf$, $\zhac\dto N(0,1)$, so the test that rejects when
$\zhac>z_{1-\alpha}$ has asymptotic level $\alpha$. $H_0$ implies $H_0'$ for every $f$, so the test
is asymptotically valid for $H_0$ whatever the critic; the dependence across windows and across
tickers on the same date is a nuisance that the date sums absorb.
\end{proposition}
\noindent\textit{Proof.} In \cref{app:proofs-dependence}.

\Cref{prop:hac} is stated for a fixed $f$, that is, for fixed training data and growing
evaluation data. Applying it with $f=\fhat$ conditionally on $\Gsig$ requires the evaluation
sequences to satisfy (D1) conditionally on $\Gsig$, in particular to be independent of $\Gsig$. If
the evaluation data depend on the training data, $\E[\bar a-\bar b\mid\Gsig]$ can be non-zero
under $H_0$, and no variance estimator corrects a bias in the mean (\cref{prop:ticker-bias}).
Versions for critics that change with the sample size would need central limit theorems uniform
over the critic class; we do not pursue them.

\begin{fact}[Mixing of the models; \citealp{carrasco2002mixing}]\label{fact:mixing}
Under $a+b<1$, with innovations that have a positive density, GARCH(1,1) is $\beta$-mixing, hence
$\alpha$-mixing, with a geometric rate; stochastic-volatility models inherit $\beta$-mixing from the
volatility process. Window-level coefficients satisfy $\alpha_W(h)\le\alpha_r((h-1)L+1)$, because a
gap of $h$ windows is a gap of at least $(h-1)L+1$ days. Finite $(2+\epsilon)$-th moments of
\GARCHt{} returns require $\nu>2+\epsilon$ and a moment condition on $(a,b)$
\citep{francq_2019_garch}.
\end{fact}

The conditional independence of training and evaluation data fails in a time series. Write
$\beta(\mathcal A,\mathcal B)=\E\sup_{C\in\mathcal B}\abs{\Prob(C\mid\mathcal A)-\Prob(C)}$ for the
absolute-regularity coefficient of two $\sigma$-fields \citep{bradley_2005_basic}.

\begin{lemma}[Decoupling]\label{lem:decoupling}
Let $U$ (training data, $\sigma$-field $\mathcal A$) and $V$ (evaluation data, $\sigma$-field
$\mathcal B$) take values in Polish spaces, and let $\tilde V$ be a copy of $V$ independent of $U$.
For every measurable $E$, $\abs{\Prob((U,V)\in E)-\Prob((U,\tilde V)\in E)}\le\beta(\mathcal A,\mathcal B)$.
\end{lemma}
\noindent\textit{Proof.} In \cref{app:proofs-dependence}.

If training and evaluation data were contiguous blocks separated by a gap of $g$ days, then
$\beta(\mathcal A,\mathcal B)\le\beta_r(g+1)$, the $\beta$-mixing coefficient of the (panel) return
process at lag $g+1$, which decays geometrically for the models of \cref{fact:mixing}. Every level
statement whose remaining hypotheses concern only the evaluation data, \cref{prop:hac} in
particular, would then hold up to an additive $\beta_r(g+1)$, because $\tilde V$ keeps the internal
dependence of $V$. The implemented split is different (\cref{sec:method-inference}): each sample's
date range is cut into $K=10$ blocks of equal length, whole blocks are assigned at random to the
two halves, the blocks interleave, and there is no gap. At every boundary between a training and
an evaluation block the relevant coefficient is the one-day coefficient, which is not small under
persistent volatility, so \cref{lem:decoupling} gives no useful bound. At the E3 sizes a \TEST{}
block holds about three window dates per ticker, so most evaluation windows lie within one or two
windows of a training block, while the volatility half-life at persistence $0.98$ is about 34
days. What the split does guarantee is that training and evaluation dates are disjoint, which
removes the common-date bias of \cref{prop:ticker-bias} below. It does not remove the bias that a
persistent common factor carries from training dates to adjacent evaluation dates
(\cref{cor:block-bias}). The residual dependence across block boundaries is controlled only
empirically: by the simulated placebo of E1, which runs the exact real-data pipeline
(\cref{sec:emp-placebo}), and by E2 part g (\cref{sec:sim-dependent}).

\begin{proposition}[Bias under ticker splits on common dates]\label{prop:ticker-bias}
Let the real windows be $X_{i\tau}=\psi(F_\tau,\varepsilon_{i\tau})$ for tickers $i$ and dates
$\tau\in\mathcal T$, where the factor path $(F_\tau)$ is strictly stationary with marginal $\mu_F$,
the $\varepsilon_{i\tau}$ are i.i.d.\ and independent of $(F_\tau)$, and $\psi$ is measurable. Let
$P^{(v)}$ be the law of $\psi(v,\varepsilon)$, so that every real window has law
$\bar P=\int P^{(v)}\mu_F(dv)$. Let the synthetic windows be mutually independent, independent of
the real data, with law $\bar Q=\bar P$ (the pooled null holds). Split the tickers into training
and evaluation sets independently of the data, each ticker contributing one window per date, and
let $\fhat$ be measurable with respect to $\Gsig\vee\sigma(F)$, for example $\Gsig$-measurable,
where $\Gsig$ is generated by the training windows (real and synthetic), the split and $\xi$.
\begin{enumerate}
  \item[(a)] Conditionally on $(\Gsig,F)$ the real evaluation windows are independent with
  $X_{i\tau}\sim P^{(F_\tau)}$, and $\E[T\mid\Gsig,F]=(\hat P_F-\bar Q)\fhat$, where
  $\hat P_F=\abs{\mathcal T}^{-1}\sum_{\tau\in\mathcal T}P^{(F_\tau)}$ is the realised mixture.
  \item[(b)] $\E T=\E[(\hat P_F-\bar Q)\fhat]$. This vanishes if $\fhat$ is independent of $F$;
  in general it is a covariance between the critic and the realised factor path and need not
  vanish.
  \item[(c)] (Idealised critic.) If $\fhat=\log(d\hat P_F/d\bar Q)$, the Bayes logit of the
  population version of the training problem given $F$ (the limit when the number of training
  tickers per date grows, so that the training half reveals $F$), then
  $\E[T\mid\Gsig,F]=\Jeff(\hat P_F,\bar Q)\in(0,\infty]$ whenever $\hat P_F\neq\bar Q$. For Gaussian
  scale mixtures, $\psi(v,\varepsilon)=v\varepsilon$ with $\varepsilon\sim N(0,I_L)$ and $v>0$, this
  holds almost surely as soon as $\mu_F$ has no atoms.
\end{enumerate}
\end{proposition}
\noindent\textit{Proof.} In \cref{app:proofs-dependence}.

With many training tickers per date, the training half reveals the realised factor path. A critic
that separates real from synthetic training windows then learns that windows typical of the
realised dates are real; the real evaluation windows share those dates and the synthetic ones do
not. The result is a bias, $\E[T\mid\Gsig]\neq0$ and typically positive, under the pooled null.
Driscoll--Kraay and other HAC variances estimate the dispersion of $T$ around $\E[T\mid\Gsig]$ and
cannot remove it; the permutation p-value, which also ignores the cross-sectional dependence, is
worse. A placebo that compares real tickers with real tickers on the same dates cannot reveal the
bias, because both samples share the realised path. The simulated placebo E1(a) runs the ticker
split with permutation p-values as a comparison row; its rejection rate at nominal 5\% is 1.000
[0.981, 1.000], and with the Driscoll--Kraay p-value still 0.990 [0.964, 0.997], as the proposition
predicts for a bias in the mean. The time-block split with the Driscoll--Kraay p-value rejects at
0.085 [0.054, 0.132] (Wilson 95\% intervals, 200 replications; \cref{tab:e1-size-sim}).

The time-block split makes training and evaluation dates disjoint, but the proof of
\cref{prop:ticker-bias}(a)--(b) does not use that the two halves share their dates. It uses only
that, given $(\Gsig,F)$, the real evaluation windows are independent with laws $P^{(F_\tau)}$ and
the synthetic ones are independent of $(\Gsig,F)$ with law $\bar Q$.

\begin{corollary}[Time-block split under a persistent factor]\label{cor:block-bias}
In the model of \cref{prop:ticker-bias}, split the dates instead of the tickers: $\mathcal T$ is
cut into disjoint sets $\mathcal T_{\trn}$ and $\mathcal T_{\evl}$, chosen independently of the
data, and every real and synthetic ticker contributes its windows on $\mathcal T_{\trn}$ to the
training half and those on $\mathcal T_{\evl}$ to the evaluation half. Let $\fhat$ be
$\Gsig$-measurable and $\hat P_{F,\evl}=\abs{\mathcal T_{\evl}}^{-1}\sum_{\tau\in\mathcal
T_{\evl}}P^{(F_\tau)}$. Then $\E[T\mid\Gsig,F]=(\hat P_{F,\evl}-\bar Q)\fhat$ and
$\E T=\E[(\hat P_{F,\evl}-\bar Q)\fhat]$. This vanishes if $(F_\tau)_{\tau\in\mathcal T_{\evl}}$ is
independent of $(F_\tau)_{\tau\in\mathcal T_{\trn}}$. In general it is a covariance between the
critic, which depends on the factor path on the training dates, and the factor path on the
evaluation dates, and it need not vanish when $(F_\tau)$ is persistent across adjacent blocks.
\end{corollary}
\begin{proof}
The real evaluation windows are functions of $(F_\tau,\varepsilon_{i\tau})_{\tau\in\mathcal
T_{\evl}}$. $\Gsig$ depends on the real data only through $(F_\tau,\varepsilon_{i\tau})_{\tau\in
\mathcal T_{\trn}}$, so the $\varepsilon_{i\tau}$ with $\tau\in\mathcal T_{\evl}$ are independent of
$(\Gsig,F)$, and the proof of \cref{prop:ticker-bias}(a)--(b) applies with $\mathcal T_{\evl}$ in
place of $\mathcal T$, under the same integrability. If $(F_\tau)_{\tau\in\mathcal T_{\evl}}$ is
independent of $(F_\tau)_{\tau\in\mathcal T_{\trn}}$, it is independent of $\Gsig$, and
$\E[P^{(F_\tau)}\fhat\mid\Gsig]=\bar P\fhat=\bar Q\fhat$ for $\tau\in\mathcal T_{\evl}$ by
stationarity.
\end{proof}

With a persistent factor the realised window law on the evaluation dates resembles the one on the
adjacent training dates, which the critic has learned to score as real. The corollary is stated,
like \cref{prop:ticker-bias}, for mutually independent synthetic windows, the situation of
univariate generators that simulate tickers independently; when the synthetic panel has a common
factor of its own, the same computation gives $(\hat P_{F,\evl}-\hat Q_{F',\evl})\fhat$ with the
realised mixture $\hat Q_{F',\evl}$ of the synthetic factor path, which need not vanish either.
\Cref{rem:small-T} compares this mechanism with the E1 measurements.

\begin{remark}[Few evaluation dates]\label{rem:small-T}
For features driven by common factors, precision is governed by the number of distinct evaluation
dates, not by tickers times windows. At $L=20$ the \TEST{} period gives 30 window dates per ticker,
about 15 of them in the evaluation half. HAC estimates from so few dates are imprecise and biased
downwards, and normal critical values then over-reject. Fixed-$b$ critical values matched to the
Bartlett kernel and bandwidth \citep{kiefer_2005_new}, or a cluster variance over time blocks with
$t_{G-1}$ critical values \citep{bester_2011_inference}, are the standard small-sample
corrections; the two pairings should not be mixed. The pipeline reports the normal-reference
Driscoll--Kraay p-value; fixed-$b$ p-values are computed as comparison rows in E1, which measures
the size of the complete pipeline (\cref{sec:emp-placebo}). There, at nominal 5\%, the \WC{} test
rejects at 0.085 with the normal reference and at 0.050 with the fixed-$b$ reference for equal
variances when the synthetic panel shares the factor structure of the real one, but at 0.270 and
0.245 when the synthetic tickers are independent (\cref{tab:e1-size-sim}). The second case is not
a small-sample variance effect alone: the mean of $\zhac$ is 0.74 there, against 0.11 in the first
case (\cref{tab:e1-sim-z,sec:sim-dependent}). \Cref{cor:block-bias} suggests a mechanism, which
is model-based and heuristic. The conditional mean $(\hat P_{F,\evl}-\bar Q)\fhat$ of $T$ is a
function of the realised factor path, and so is the Driscoll--Kraay standard error, through the
date-level variance of the real critic values. When the synthetic tickers are independent, their
date sums carry no common factor, the standard error is driven by the real path alone, and a
conditional bias that is small on average can still shift $\zhac$ if it is dependent on that
standard error. A synthetic factor of its own adds variation to the standard error that is
independent of the real path, which weakens this dependence.
\Cref{tab:manual-e1-sim-stat} is consistent with this reading. The mean of $T$ for the \WC{} test
is small and nearly the same in the two cells (0.026 and 0.029, with Monte Carlo standard errors
0.015 and 0.014), but with independent synthetic tickers $T$ and its standard error have
Pearson correlation $-0.64$, against 0.22 in the joint cell, so that positive values of $T$ come with
small standard errors; the \LRC{} test shows the same sign pattern, and the rejections concentrate
in the replications with a small standard error (\cref{tab:e1-sim-se}). The E1 panel lies outside the
model of the corollary, because its idiosyncratic and synthetic GARCH components are also
serially dependent within a ticker, and the corollary does not quantify the effect.
\end{remark}

\begin{table}[tbp]
  \centering
  \caption{Statistic and Driscoll--Kraay standard error of the absolute learned tests in the E1(a)
    placebo (time-block split, 200 replications per cell). joint: the synthetic panel is an
    independent draw of the whole factor panel; indep.: independent synthetic tickers. Mean of
    $T$ with its Monte Carlo standard error, standard deviation of $T$, share of replications with
    $T>0$, median Driscoll--Kraay standard error of $T$, correlation of $T$ with that standard
    error over the replications, and mean of $\zhac$ (as in \cref{tab:e1-sim-z}). Derived from
    \texttt{results/e1/e1\_sim\_raw.csv} by \texttt{paper/scripts/e1\_sim\_stat.py}.}
  \label{tab:manual-e1-sim-stat}
  \small
  \fittable{manual_e1_sim_stat}
\end{table}

\begin{remark}[Wild bootstrap for the kernel tests]\label{rem:wild}
MMD and \SigMMD{} have no train/evaluation split. Under dependence they use the wild bootstrap of
\citet{chwialkowski2014wild} for the MMD V-statistic, with a Gaussian AR(1) multiplier process on
the date grid that is shared by all windows with the same date rank, in both samples, and a
plug-in correlation length (\cref{app:impl}). \citet{chwialkowski2014wild} prove validity for
weakly dependent stationary processes with a given correlation length; the date-shared multipliers
and the plug-in length are our choices and are not covered by their theorem. We give no proof;
their size is measured in E1 and in E2 part g.
\end{remark}

% ---------------------------------------------------------------------------------------------
\subsection{What the verdict controls}\label{sec:theory-verdict}

\begin{definition}[Power check]\label{def:power-check}
A power check is a pre-specified pair $(\Cdesign,\vartheta^*)$. $\Cdesign$ is a reference
simulation design that reproduces the empirical one: the sizes of $\Rpast$, $\Rnew$ and $G$, the
window length, the split, the critic settings, the inference, the seed combination and, where
possible, the dependence. $\vartheta^*$ is a pre-specified deviation of the generator. The
boundary generator $G^0$ is an independent draw with the law of the past windows. We use two
designs. D1 is a per-ticker \GARCHt{} panel fitted on \TRAIN: $\Rpast^*$ is a \TRAIN-length path
started from stationarity, $\Rnew^*$ a \TEST-length path started from the fitted end-of-\TRAIN{}
state, $G^0$ a \TEST-length path started from stationarity, and the alternative $G^*$ is $1.2$
times an independent \TEST-length path started from the end-of-\TRAIN{} state. In D2 the past
windows are the real \TRAIN{} realised volatilities times fresh Gaussian noise, $\Rnew$ is the
real \TEST{} sample, $G^0$ has the pooled law of the past windows, and $G^*=1.2\,G^0$. In D2 the
alternative differs from the boundary only by the volatility factor; in D1 it also differs by its
initial state, so the D1 power is the power against a compound deviation, not against the
volatility factor alone. The power $\beta(\vartheta^*)$ is the probability that the verdict rule
rejects under $G^*$, estimated from $R$ replications by $\hat\beta$, with one-sided
Clopper--Pearson lower bound $\betaL$ \citep{clopper_1934_use}. The same design with $G^0$ in place
of $G^*$ gives the boundary size $\abd$, estimated by $\hat\alpha_{\mathrm{bd}}$. Since $\abd$ can
exceed $\alpha$, $\hat\beta$ mixes the discrimination between $G^*$ and $G^0$ with the
miscalibration of the rule.
The size-adjusted power $\beta_{\mathrm{adj}}(\vartheta^*)=\Prob(p^{\mathrm{alt}}\le c_\alpha)$,
where $p^{\mathrm{alt}}$ is the rule's $\pmed$ under $G^*$ and $c_\alpha$ the $\alpha$-quantile of
its law $p^{\mathrm{bd}}$ under $G^0$, measures the discrimination alone. It is estimated with
$\hat c_\alpha$, the $\lfloor\alpha R\rfloor$-th smallest of the $R$ boundary values, at which the
empirical boundary size is $\lfloor\alpha R\rfloor/R\le\alpha$. We recommend reporting
$\hat\beta_{\mathrm{adj}}$ together with $\hat\alpha_{\mathrm{bd}}$ and its Clopper--Pearson
interval.
\end{definition}

\begin{definition}[Verdict]\label{def:verdict}
Let $\alpha=0.05$, $\bmin=0.8$ and let $\pmed$ be the anchored p-value combined over the synthetic
seeds (\cref{rem:combine-seeds}); as implemented, each single-seed p-value is the Driscoll--Kraay
p-value $p_s=1-\Phi(\zhac)$ of $\Danc$. The verdict \eqref{eq:verdict} is $V=\BAD$ if
$\pmed\le\alpha$; $V=\GOOD$ if $\pmed>\alpha$ and $\hat\beta\ge\bmin$ under every design (D1 and D2); and
$V=\INCONCL$ otherwise. The stationary bootstrap, for which the anchored test is invalid
(\cref{rem:anchored-null}), receives \INVALID. A stricter variant replaces $\hat\beta$ by $\betaL$.
\end{definition}

\begin{proposition}[What the verdict controls]\label{prop:verdict}
(a) If each single-seed anchored p-value $p_s$ satisfies $\Prob(p_s\le u)\le u$ for all
$u\in(0,1)$, then $\Prob(V=\BAD)\le\alpha$. For the permutation version of the anchored p-value
this premise holds in finite samples at $\QG=\Ppast$ under \cref{ass:anc} (\cref{prop:anchored}(a)).
The rule as implemented uses the Driscoll--Kraay p-values (\cref{def:verdict}), for which the
premise holds only in the limit, also when the windows are independent: at $\QG=\Ppast$ under the
conditions of \cref{prop:anchored}(c), where $\delta=0$ by \cref{prop:anchored-target}(a),
$\limsup\Prob(V=\BAD)\le\alpha$. On real data, where the past is non-stationary
(\cref{rem:nonstat-past}), the premise does not hold; there the boundary size $\abd$ of the rule
is estimated (\cref{def:power-check}), and a \BAD{} verdict has a real-data error interpretation
only when the Clopper--Pearson interval of $\hat\alpha_{\mathrm{bd}}$ covers $\alpha$ or lies below
it.
(b) If the rejection probability of the rule at the true configuration is at least
$\beta(\vartheta^*)$, then $\Prob(V=\GOOD)\le\Prob(\pmed>\alpha)\le1-\beta(\vartheta^*)$, which is
at most $1-\bmin$ up to the Monte Carlo error of $\hat\beta$.
\end{proposition}
\begin{proof}
(a) $\{V=\BAD\}=\{\pmed\le\alpha\}$; apply \cref{rem:combine-seeds}. Under
\cref{prop:anchored}(c), its Markov bound $\Prob(\pmed\le\alpha)\le(2/S)\sum_s\Prob(p_s\le\alpha/2)$
holds for every sample size, and $\limsup\Prob(p_s\le\alpha/2)\le\alpha/2$ for each seed follows
from the conditional limit by dominated convergence.
(b) $\{V=\GOOD\}\subset\{\pmed>\alpha\}$.
\end{proof}

That the Driscoll--Kraay version is valid only in the limit is visible at the sizes used. The
single-seed Driscoll--Kraay anchored p-value rejects at the boundary in 0.09 [0.04, 0.16] of the D1
replications (Clopper--Pearson; \cref{tab:e3-power}) and in 0.080 [0.050, 0.126] of the
replications of the joint cell of E1 (Wilson; \cref{tab:e1-size-sim}), two stationary simulated
panels; in D1 the seed combination absorbs the excess, with a boundary size of the rule of 0.00
[0.00, 0.04]. The \BAD{} condition at the end of (a) is a reporting diagnostic, not a change to the
pre-specified rule; \cref{rem:calibrated-bad} gives the calibration that replaces the nominal level
when it fails.

Part (b) is only as strong as its premise, which cannot be verified: the true deviation must be at
least as detectable as $\vartheta^*$ under $\Cdesign$. \GOOD{} therefore means ``no detectable
deviation of the pre-specified size and type'', never ``equal to the truth''. Our $\vartheta^*$ is
an upward volatility error, which in D2 points away from the calmer \TEST{} period. Power against
downward volatility errors or errors in other features is not covered (a factor 0.8 was not
run), and under a non-stationary past a non-rejection of the anchored null can reflect a generator
bias aligned with the regime change (\cref{rem:anchored-null}). The D2 boundary size, in turn,
is that of a boundary generator with Gaussian windows whose scales are drawn independently across
tickers and windows (\cref{sec:method-decision}). D1 is a design-stage
calculation, like a sample-size calculation. D2 is a power analysis conditional on the realised
$\Rnew$ and on the \TRAIN{} volatility path: it uses the evaluation data, though not the generators,
so \GOOD{} against \INCONCL{} depends on $\Rnew$ through $\pmed$ and through $\hat\beta$. Part (b)
is unaffected, because $\{V=\GOOD\}\subset\{\pmed>\alpha\}$. With $R=100$ replications the standard
error of $\hat\beta$ near $0.8$ is $0.04$, which motivates the Clopper--Pearson variant.

\begin{remark}[Calibrating \BAD{} against the boundary replications]\label{rem:calibrated-bad}
Let $p^{\mathrm{bd}}_1,\dots,p^{\mathrm{bd}}_R$ be the rule's $\pmed$ in the $R$ boundary
replications of a design and $p^{\mathrm{obs}}$ the observed $\pmed$ of a generator. The Monte
Carlo p-value
\[
  p_{\mathrm{MC}}=\frac{1+\#\{r:p^{\mathrm{bd}}_r\le p^{\mathrm{obs}}\}}{R+1}
\]
satisfies $\Prob(p_{\mathrm{MC}}\le u)\le u$ for every $u\in(0,1)$ if, under the null,
$(p^{\mathrm{obs}},p^{\mathrm{bd}}_1,\dots,p^{\mathrm{bd}}_R)$ is exchangeable: apply
\cref{lem:rank} to the negated values. This is the Monte Carlo test principle
\citep{dufour_2006_monte}, with ties resolved conservatively, and its validity rests entirely on
the exchangeability. D2 keeps the real $\Rnew$, as the observed value does, but provides
exchangeability only approximately. Its past windows keep the real \TRAIN{} volatility path but
not the real within-window law (Gaussian noise, so no conditional tails, within-window clustering
or leverage) nor the same-date return correlation across tickers; its boundary generator $G^0$
has the pooled \TRAIN{} window law but not the structure of a fitted model; and on real data the
standard error of $\Danc$ is 0.59 to 0.79 times its D2 value (\cref{app:emp-derived}). The Monte
Carlo p-values of \cref{sec:emp-verdicts} are therefore approximate calibrations, more
informative than the nominal level when $\hat\alpha_{\mathrm{bd}}>\alpha$, but not exact. The
size-adjusted power of \cref{def:power-check} separates discrimination from miscalibration
(\cref{tab:manual-theory-calibration}). At the main split the D2 power of the rule is 0.45 at the
nominal threshold, with boundary size 0.08, and still 0.45 at $\hat c_\alpha=0.025$, where the
boundary size is 0.05; in all 100 replications the alternative $\pmed$ lies below the boundary
$\pmed$ of the same replication. At the alternative split of E5 the D2 power and boundary size are
both 0.60 at the nominal threshold, and the size-adjusted power is still 0.60 at
$\hat c_\alpha=3.8\times10^{-5}$ (boundary size 0.04). There the rule still separates the boundary
generator from one with 20\% too much volatility, but the Clopper--Pearson interval of
$\hat\alpha_{\mathrm{bd}}$ lies above $\alpha$, so a nominal \BAD{} verdict would have no
real-data error interpretation.
\end{remark}

\begin{table}[tbp]
  \centering
  \caption{Calibration of the verdict rule in the power-check designs: nominal boundary size and
    power at $\alpha=0.05$, estimated $\alpha$-quantile $\hat c_\alpha$ of the boundary $\pmed$
    (the $\lfloor\alpha R\rfloor$-th smallest of the $R$ boundary values), boundary size and
    size-adjusted power at $\hat c_\alpha$ (\cref{def:power-check}), and the number of replications
    in which the alternative $\pmed$ is below the boundary $\pmed$ of the same replication. Main
    split: E3 (\cref{tab:e3-power}); alternative split: E5 (\cref{tab:manual-e5-split}). $R$
    replications, $S$ synthetic seeds. Derived from \texttt{results/e3/e3\_power\_raw.csv},
    \texttt{e3\_realvol\_raw.csv} and \texttt{results/e5/e5\_split\_power\_raw.csv} by
    \texttt{paper/scripts/theory\_calibration.py}.}
  \label{tab:manual-theory-calibration}
  \small
  \fittable{manual_theory_calibration}
\end{table}

\begin{remark}[Equivalence-test alternative]\label{rem:equivalence}
A \GOOD{} verdict with a formal error guarantee would follow from a non-inferiority test
\citep{schuirmann_1987_comparison}: test $\delta\ge\varepsilon_0$ against $\delta<\varepsilon_0$ and
declare \GOOD{} if the upper $(1-\alpha)$ confidence bound for $\delta$ lies below a margin
$\varepsilon_0$, for instance the expected anchored discrepancy at $\vartheta^*$ in the power-check
design. The guarantee would again be conditional on the trained critic, and the margin transfers
across configurations only to the extent that critic training behaves similarly (heuristic). We do
not use it for the verdict.
\end{remark}

\begin{remark}[Per-ticker tests and multiple testing]\label{rem:bh}
The per-ticker tests of E3 ($L=5$) give p-values $p_1,\dots,p_N$, calibrated against a pooled
\GARCHt{} null (\cref{sec:emp-perticker}). The Benjamini--Hochberg procedure at level $q$ controls
the false discovery rate of \BAD{} verdicts at $qN_0/N$ under independence or positive regression
dependence \citep{benjamini1995controlling,benjamini2001control}. Per-ticker statistics are
plausibly positively dependent through the market factor, which we do not prove, so the
Benjamini--Yekutieli correction is reported as a robustness check and Storey's $\hat\pi_0$
\citep{storey_2002_direct} as a summary. \GOOD{} is a non-rejection, so no false-discovery
statement applies to it, and with about 60 evaluation windows per ticker few per-ticker \GOOD{}
verdicts can be expected.
\end{remark}

% ---------------------------------------------------------------------------------------------
\subsection{Window distances and portfolio utility}\label{sec:theory-markowitz}

The Markowitz experiment E4 asks whether test results predict the utility of portfolios built
from synthetic data. The theory of \cref{sec:mk-theory} (proofs in \cref{app:proofs-markowitz})
predicts a link that is real but not monotone. For the unconstrained global minimum-variance
portfolio $\wgmv(\hat\Sigma)=\hat\Sigma^{-1}\one/(\one'\hat\Sigma^{-1}\one)$, evaluated with the
\TEST{} covariance $\Sigma$, the variance loss is a squared Mahalanobis distance to the ex-post
optimum (\cref{lem:quad-loss}), the weights depend on $\hat\Sigma$ only up to scale, and
\begin{equation}\label{eq:kantorovich-preview}
  1\le\frac{V(\wgmv(\hat\Sigma))}{V(\wgmv(\Sigma))}\le\frac{(1+\kappa)^2}{4\kappa},
  \qquad \kappa=\frac{\lambda_{\max}(\Sigma^{-1/2}\hat\Sigma\Sigma^{-1/2})}
  {\lambda_{\min}(\Sigma^{-1/2}\hat\Sigma\Sigma^{-1/2})},
\end{equation}
a sharp bound that is quadratic in the relative covariance error (\cref{prop:gmv-sensitivity}).
In the large-sample limit the GMV utility of a generator therefore depends on it only through its
covariance matrix, up to scale (\cref{cor:gmv-drivers}); tails, volatility clustering and jumps,
which the critic detects, matter only through that matrix. The two criteria can disagree in both
directions. Here $\kappa$ is the ratio of the extreme eigenvalues of $\Sigma^{-1}\hat\Sigma$,
written $\cond(\Sigma^{-1}\hat\Sigma)$ in \cref{sec:mk-theory}. A generator whose covariance is
$c\Sigma$ with $c\neq1$ is rejected by a consistent
critic test and loses nothing for GMV (\cref{ex:scale-blind}); a covariance error in a
low-variance spread direction attains the bound \eqref{eq:kantorovich-preview} while $\W$ between
the window laws tends to zero (\cref{ex:spread}); univariate tests cannot see correlations
(\cref{ex:corr-blind}); and a drift error of one asset equal to one \TRAIN{} standard error
dominates mean-variance utility, yet by \cref{prop:oracle-bound}(b) no level-5\% two-sample test
that uses the 615 \TEST{} days of that asset detects it with power above $0.17$
(\cref{ex:drift}). Finally, under convex
losses no portfolio rule trained on synthetic data beats the best rule based on the real past
(\cref{prop:no-free-lunch}); synthetic data can only help as a regulariser. These results motivate
the design of E4 (\cref{sec:mk-link}): an association between critic distance and GMV loss is
plausible through covariance accuracy, little association is expected for mean-variance
portfolios, which are driven by drift, and univariate verdicts carry no information about
correlations. Because the real correlations themselves changed between \TRAIN{} and \TEST{}
(\cref{tab:e4-correlation}), fidelity to the past and usefulness for the future are different
targets.

% ---------------------------------------------------------------------------------------------
\begin{table}[tbp]
  \centering
  \caption{Status of the theoretical results. ``Proved'': full proof in \cref{app:proofs} or in
  the text; ``sketch'': key steps only; ``heuristic'': argument, not a theorem.}
  \label{tab:theory-status}
  \small
  \begin{tabular}{@{}p{0.43\linewidth}p{0.17\linewidth}p{0.34\linewidth}@{}}
    \toprule
    Result & Status & Key assumptions \\
    \midrule
    \Cref{prop:window-ident,prop:markov-pairs}: what windows identify & proved & stationarity; both processes Markov \\
    \Cref{prop:unbiased}: $\E[T\mid\Gsig]=\Dcrit(\fhat)\le\Lip(\ghat)\W(P_A,Q_A)$ & proved & evaluation data independent of training \\
    \Cref{prop:exact-level}: exact level of the permutation test & proved & conditional exchangeability (A1$'$) \\
    \Cref{prop:consistency,cor:w1-limit}: consistency; $T\to\W$ & proved & (C1)--(C5), incl.\ hard Lipschitz constraint \\
    \Cref{prop:bayes-logit,cor:jeffreys}: Bayes logit, Jeffreys & proved & perfect discriminator \\
    \Cref{prop:oracle-bound}: oracle bound & proved & densities; level at the null point \\
    \Cref{prop:local-power}: local power factor & proved (fixed $f$); sketch (learned) & differentiability in quadratic mean \\
    \Cref{rem:lr-vs-w1}: \LRC{} vs \WC{} & heuristic & -- \\
    \Cref{prop:anchored}(a): anchored test exact at $\QG=\Ppast$ & proved & \cref{ass:anc} \\
    \Cref{prop:anchored}(b)--(c): weaker null; HAC & proved (standard results) & i.i.d.\ or \cref{ass:mixing}, given $\Gsig$ \\
    \Cref{prop:anchored-target}: meaning of $\delta$ & proved & idealised critics in (b), (c) \\
    \Cref{rem:nonstat-past}: non-stationary past & heuristic & -- \\
    \Cref{prop:paired-boot}: paired bootstrap & sketch & \cref{ass:rel} \\
    \Cref{prop:boundary-size}: two-critic size $>\alpha$ at boundary & proved (given \cref{prop:paired-boot}) & \cref{ass:rel}; exchangeable training \\
    \Cref{rem:naive-perm}: naive permutation distortion & heuristic & -- \\
    \Cref{prop:hac}: HAC test & proved (standard results) & \cref{ass:mixing}; fixed critic \\
    \Cref{lem:decoupling}: decoupling & proved & no useful bound without a gap \\
    \Cref{prop:ticker-bias}: ticker-split bias & proved & factor model \\
    \Cref{cor:block-bias}: time-block bias under a persistent factor & proved & factor model; independent synthetic windows \\
    \Cref{rem:small-T}: mechanism of the E1 oversize & heuristic & -- \\
    \Cref{prop:verdict}: verdict & proved under premises & valid $p_s$ (asymptotic only for the Driscoll--Kraay $p_s$ used); unverifiable power premise \\
    \Cref{rem:calibrated-bad}: Monte Carlo calibration of \BAD{} & proved (Monte Carlo test) & exchangeable observed and boundary $\pmed$ (approximate in D2) \\
    \Cref{sec:mk-theory}: GMV and utility results & proved; parts heuristic & positive definite matrices; convex loss \\
    \bottomrule
  \end{tabular}
\end{table}
